III) \emph{Caso en que $A$ es un cuerpo, $B$ un cuerpo que es una
extensión de tipo finito de $A$, y $A$ está algebraicamente cerrado
en $B$.} --- Entonces $C=A$, y $A'$ es una $A$-álgebra noetheriana
geométricamente normal, por lo que es un producto directo de
dominios enteros integralmente cerrados $A'_i$, siendo los cuerpos de
fracciones $K'_i$ de los $A'_i$ extensiones \emph{separables} de $A$ (4.6.1). La
$A'$-álgebra $B'$ es, por tanto, el producto directo de los
$B\otimes_A A'_i=B'_i$, y es posible limitarse al caso en que $A'$ es
un dominio entero y su cuerpo de fracciones $K'$ una extensión
separable de $A$. Como $B$ es un $A$-módulo plano,
$B\otimes_A A'=B'$ se identifica con un subanillo de
$B\otimes_A K'$. Puesto que $K'$ es una extensión separable de $A$,
$B\otimes_A K'$ es reducido (4.3.7); y como $B$ es además una
extensión primaria de $A$, $B\otimes_A K'$ es un dominio entero
(4.3.2). Por tanto, $B'$ también lo es. Sea $L'$ el cuerpo de
fracciones de $B\otimes_A K'$, que es asimismo el de $B'$; es
evidentemente el cuerpo compuesto $B(K')$ de sus subcuerpos $B$ y
$K'$. Como $A$ está algebraicamente cerrado en $B$, $K'$ es una
extensión separable de $A$, y $B$ y $K'$ son linealmente disjuntos
sobre $A$, el cuerpo $K'$ está algebraicamente cerrado en
$B(K')=L'$ (Bourbaki, \emph{Alg.}, cap. V, § 9, ejerc. 2); \emph{a fortiori}, $A'$, al ser integralmente cerrado, está integralmente
cerrado en $L'$ y, por tanto, en $B'$, lo que demuestra la
proposición en este caso.

